\section*{Limitations}
\label{sec:limitations}

While we believe that \method{} offers a valuable contribution to context parameterization and utilization for continually updated knowledge, several limitations should be acknowledged. First, \method{} is built entirely atop D2L, whose hypernetwork compresses a context into a compact set of LoRA parameters; this parameterization is inherently lossy, and long or information-dense contexts may not be faithfully preserved during encoding, which in turn upper-bounds the overall performance of \method{} by the fidelity and capacity of the underlying parameterization method. Second, \method{} is deliberately training-free, a design choice that preserves its flexibility, low memory footprint, and low latency, but that also confines its interventions to parameter construction and decoding time, without directly enhancing the hypernetwork's intrinsic ability to discriminate valid updates from invalidated information. Our future work will therefore focus on jointly improving the fidelity of context parameterization and equipping the hypernetwork itself with update-aware, training-based mechanisms, aiming to provide continued valuable insights for the development of truly reliable context-adaptive systems.